\section{Markowitz and Its Discontents}

Markowitz~\cite{markowitz1952} framed portfolio selection as a quadratic program: choose weights $\mathbf{w} \in \mathbb{R}^p$ to
\begin{equation}
\min_{\mathbf{w}}\; \mathbf{w}^{\!\top}\boldsymbol{\Sigma}\mathbf{w} \quad \text{s.t.} \quad \mathbf{w}^{\!\top}\boldsymbol{\mu} \geq \mu_0,\;\; \mathbf{w}^{\!\top}\mathbf{1} = 1.
\label{eq:markowitz}
\end{equation}
The solution delivers the familiar efficient frontier --- a parabola in mean-variance space. Its weaknesses are equally familiar: $\boldsymbol{\mu}$ and $\boldsymbol{\Sigma}$ are notoriously difficult to estimate (the ``Markowitz error-maximization'' problem~\cite{michaud1989}), and the Gaussian assumption ignores tail dependence.

Subsequent decades produced a rich literature on robust optimization~\cite{fabozzi2007}, shrinkage estimation~\cite{ledoit2004}, and Black--Litterman blending~\cite{black1992}. These refinements improve \textit{static} estimates but do not address the core problem: the data-generating process itself is non-stationary.
